\subsection{Slope Trick}

\begin{itemize}
    \item A function is slope-trick-able if it satisfies 3 conditions:


\begin{enumerate}
    \item It is continuous.
    \item It can be divided into multiple sections, where each section is a linear function with an integer slope.
    \item It is convex/concave.
\end{enumerate}

\item A slope-trick-able function can be represented by a linear function of the rightmost section and a multiset \textbf{S} containing all the slope changing points where the slope changes by 1.

\item If $f(x)$ and $g(x)$ are slope-trick-able, then so is $f(x) + g(x)$ (merge the multiset and sum the linear functions).

\item Example: given an array, each operation you are allowed to increase or decrease an element's value by 1. Find the minimum number of operations to make the array non decreasing.

\item $f_i(x) = $ minimum number of operations to make the first $i$ elements of the array non-decreasing, with the condition that $a_i \leq x$ is slope-trick-able.

\item  The dual of the problem is: given an array with the price of some stock by day, each day we can either buy or sell one unity of stock or do nothing. What is the largest possible profit?

\end{itemize}
